\section{The next breakthrough may be the ability to build it}
\label{sec:opening}

Imagine a laboratory that could test an important new material, except that the necessary measurement cannot yet be made reliably. An algorithm can propose the experiment. A designer can draw the apparatus. A manufacturer can make its parts. Yet the useful capability may remain absent: the parts do not interface, calibration is missing, the reference sample is unavailable, or no one can fund the sequence before an accepted result exists. This is an illustrative situation, not an experiment reported in this paper.

The consequential construction is the complete means of performing the measurement. Its output is an ability: someone can now make a distinction, produce a result, or undertake useful work that was previously inaccessible under the stated conditions. The construction can include familiar components. Its significance lies in what they make possible together.

Now consider a second laboratory problem. Keeping the first apparatus may help. Keeping its design principles may help a different apparatus. Learning a better way to identify missing measurements, construct fixtures, diagnose failures, and qualify new designs would be a further achievement. The object that improves has changed from an instrument to the process of making instruments.

\begin{thesisbox}
\textbf{Learning to build is improving the ability to turn a worthwhile possibility into working, accepted means---and to make fresh construction better through what was learned.}
\end{thesisbox}

This is the connective tissue between an ambitious idea and an expanding productive civilization. A useful instrument can improve science. Better science can improve processes. A qualified process can make further instruments, energy systems, or production facilities affordable. Their operation can supply resources and knowledge for another construction. Every connection is a claim to be tested; the sequence is a research and enterprise objective, not an inevitable trajectory.

\subsection{The unit of progress is a changed ability}

An answer is often exactly what is needed. So is an existing tool, a purchased service, or a simple process change. Construction becomes relevant when an important outcome remains blocked because the required means are absent, incompatible, unreliable, or unaffordable. The right response need not be a larger system. It may remove a requirement that belonged only to the old method, or replace an intricate solution with a mature alternative.

The phrase \emph{what makes new possibilities real} therefore has an operational meaning. State an accepted outcome; identify the initial conditions; construct a permitted route; execute and check it at the claimed evidential level. A digital artifact can be a real computational construction. A simulated factory is a model of a factory. A proposed sale is a hypothesis about a relationship. These distinctions preserve the meaning of each achievement.

We use \emph{worthwhile} to refer to a beneficiary's authorized purpose, with relevant costs, risks, and obligations included. It is not a license for the builder to invent beneficiaries' preferences. The choice of what to make is part of the problem, but acceptance belongs to the people and institutions whose commitments it requires.

\subsection{Three foundations, one next question}

\emph{Transformational Intelligence} distinguishes an intended purpose, a formulation for pursuing it, and the instrument used to judge success \citep{boucherTI}. A change of representation, observation, dependency, or coordination can remove a restriction that more search inside the old formulation cannot overcome. The change still needs an executable realization and a bridge back to the purpose.

\emph{Invent the Economy} makes the productive arrangement itself an object of invention \citep{boucherEconomy}. Technology, participation, prices, risk allocation, and the timing of funding may have to change together. Its separate claims about technical possibility, modeled viability, accepted arrangements, exchange, and repeated adoption remain essential.

\emph{ALPHA GENESIS FLYWHEEL $\Omega$} connects frontier developments to new valuable possibilities and to improvement of the machinery that discovers them \citep{boucherGenesis}. It already describes complementary capabilities, a construction grammar that must remain open to revision, and causal tests of successor benefit.

The present paper puts a further question at the center:
\begin{center}
\large\color{Plum}\textbf{How does an intelligence become better at constructing\\the complete means of a new possibility?}
\end{center}
Its contribution is an operational synthesis and a sharper unit of evaluation. We define a complete construction witness, distinguish installed and constructible capability, make bootstrap and ongoing obligations explicit, and propose transfer tests in which the donor's finished artifact is unavailable. These ideas develop the trilogy; they do not claim to originate systems engineering, program synthesis, economic complementarity, or learning from experience.

\subsection{A bridge between established fields}

The engineering side has a mature foundation. NASA's systems engineering guidance treats interactions, integration environments, verification, validation, support, and transition as part of realizing a system; its discussion includes the tooling needed to assemble and test other products \citep{nasaSE}. The learning target here is the discovery and improvement of such complete construction processes across new missions.

Computational creativity distinguishes combination, exploration, and transformation of conceptual spaces \citep{boden1998}. DreamCoder learns symbolic abstractions and search guidance, while LAPS examines language-guided learning of program libraries and search heuristics \citep{ellis2020,wong2021}. POWERPLAY searches jointly over new tasks and modifications of a solver \citep{schmidhuber2011}. AlphaEvolve provides a concrete precedent for evaluator-guided program improvement in scientific and computational applications \citep{novikov2025}. These systems motivate parts of the agenda. Their results do not supply the untested physical, contractual, or economic bridges required here.

Complementarity explains why a coherent set of technical and organizational changes can matter more than isolated upgrades \citep{milgrom1990}. Endogenous-growth theory supplies a precedent for treating investments in knowledge as part of economic production \citep{romer1990}. Neither precedent makes an individual proposed venture profitable or implies a universal rate of expansion.

\begin{table}[htbp]
\centering\small
\begin{tabularx}{\linewidth}{@{}P{1.31in}YY@{}}
\toprule
\textbf{Research object} & \textbf{What it helps explain} & \textbf{Question emphasized here}\\
\midrule
Systems engineering & Integration and realization of a useful whole. & Can experience improve the construction of a different whole?\\
Program/library learning & Reusable abstractions and better program search. & Does the gain survive complete realization and equal-content controls?\\
Transformational intelligence & Changing a restrictive problem-solving setup. & What must actually be built to make the change operational?\\
Economic invention & Joint capability and voluntary arrangement. & Can the whole sequence be started, delivered, and sustained?\\
GENESIS & Repeated opportunity discovery and successor effects. & Which part of the builder improved, and what fresh evidence shows it?\\
\bottomrule
\end{tabularx}
\caption{A connection among established ideas and the three foundations, rather than a claim that their mechanisms are new.}
\end{table}

\section{Complete means: an executable whole}
\label{sec:means}

\subsection{Purpose first; the method remains editable}

Let a construction contract be
\begin{equation}
\I=(\mathcal Y,\mathcal H,\mathcal U,B,T,V),
\label{eq:contract}
\end{equation}
where $\mathcal Y$ describes the intended outcome, $\mathcal H$ the hard requirements, $\mathcal U$ the authority boundary, $B$ the initial resource envelope, $T$ the horizon, and $V$ the acceptance procedure. The environment and uncertainties are recorded separately. A finite model can stipulate them; a natural application must investigate whether its model is adequate.

The construction may change the method, representations, interfaces, or arrangements allowed by this contract. It cannot retrospectively improve a result by dropping an inconvenient requirement. If a better interpretation of the purpose is proposed, authorization and a new comparison make the change explicit. An acceptance instrument may also be defective. Its repair requires evidence and a common re-evaluation of candidates, not unilateral self-approval.

\begin{definition}[Complete means, relative to a contract]
Complete means are an executable combination of knowledge, instruments, processes, resources, and arrangements that can be initialized, operated, checked, and sustained sufficiently to deliver the accepted outcome within the declared scope and horizon. Necessary external dependencies, permissions, and remaining obligations are part of that declaration.
\end{definition}

\emph{Complete} does not mean self-sufficient in every respect. A software instrument may legitimately depend on an operating system, electricity, and a licensed solver. A laboratory may use an external calibration service. The question is whether a material dependency has a supported place in the route, not whether one organization can reproduce the entire world.

\begin{figure}[htbp]
\centering
\begin{tikzpicture}[x=1in,y=1in,every node/.style={font=\small,align=center},box/.style={draw=Rule,fill=Ivory,rounded corners=3pt,text width=1.76in,minimum height=.62in,inner sep=5pt},arr/.style={-{Latex[length=2.1mm]},line width=.8pt,color=Teal}]
\node[box] (k) at (0,1.3) {\textbf{KNOWLEDGE}\\Models, distinctions, evidence};
\node[box] (i) at (2.25,1.3) {\textbf{INSTRUMENTS}\\Ways to observe and act};
\node[box] (p) at (4.5,1.3) {\textbf{PROCESSES}\\Executable coordination};
\node[box] (r) at (1.1,0) {\textbf{RESOURCES}\\Time, capacity, cash, materials};
\node[box] (a) at (3.4,0) {\textbf{ARRANGEMENTS}\\Rights, terms, participation};
\node[draw=Gold,fill=Gold!7,rounded corners=3pt,text width=5.6in,minimum height=.62in,inner sep=6pt] (out) at (2.25,-1.35) {\textbf{A WORKING ABILITY}\\Initialized, integrated, accepted, and supported for its stated use};
\draw[arr] (k.south)--++(0,-.2)-|(out.north west);
\draw[arr] (p.south)--++(0,-.2)-|(out.north east);
\draw[arr] (i.south)--(out.north);
\draw[arr] (r.south)--(out.north -| r.south);
\draw[arr] (a.south)--(out.north -| a.south);
\end{tikzpicture}
\caption{The components enable an outcome together. Lines mean required integration, not independent sufficient causes. A missing or incompatible element can defeat the whole.}
\end{figure}

\subsection{The witness that connects the parts}

For a starting state $x_0$, write a proposed construction witness as
\begin{equation}
w=(\pi,D,\Sigma,A,E,O).
\label{eq:witness}
\end{equation}
Here $\pi$ is the executable construction and operating policy; $D$ records dependencies and interface meanings; $\Sigma$ schedules resource use; $A$ specifies required authority and arrangements; $E$ links claims to evidence; and $O$ records continuing obligations. A policy chooses from available observations. It cannot choose a branch using a hidden environmental fact that no allowed observation reveals.

A \emph{prospective witness} says how a modeled route could work. A \emph{realized witness} additionally records the actual execution, required acceptance, and resulting obligations. A forecast of consent belongs to the former. Only a corresponding authorized action belongs to the latter.

Five questions turn a parts list into a construction:
\begin{enumerate}
\item \textbf{Can it start?} Every initial prerequisite is available or reached by a feasible preceding action. Mutual promises to create each other do not start a process.
\item \textbf{Do the parts compose?} Interfaces have compatible meanings, units, versions, and operating conditions. A shared resource is not allocated twice.
\item \textbf{Can every step occur?} Time, liquidity, capacity, and authority hold at every prefix of execution, not only in the final totals.
\item \textbf{Does the intended outcome follow?} Acceptance checks the useful whole and its adverse cases, including failures shared by several components.
\item \textbf{What remains afterward?} Support, maintenance, debts, refunds, decommissioning, and retained capabilities have explicit owners and costs.
\end{enumerate}

These are substantive properties of the route. They need not become five separate committees, scores, or mandatory software modules. A small problem may be settled by one clear calculation and an executable test. A high-consequence physical system requires evidence proportionate to its consequences.

\subsection{Completeness has a boundary}

For an explicitly modeled environment $\theta$, let $\operatorname{Valid}(w;\I,x_0,\theta)$ mean that the policy's execution respects the contract, its resource prefixes, and its acceptance conditions. Robust model-relative completeness requires this for every $\theta$ in a declared set $\Theta$. A probabilistic version requires a stated reliability target under a justified distribution instead.

Neither version establishes that the modeled set contains every relevant real failure. An omitted disturbance, false interface assumption, or forged source can invalidate the conclusion. The witness must therefore distinguish a proved implication, a measured property, an assumption, and an unresolved dependency. Unknown conditions can motivate a new experiment or a construction that makes them irrelevant. They cannot become true because the narrative needs them.

The practical aim is \emph{enough justified completeness for the useful outcome}, with the residual boundary visible. Requiring omniscience would prevent action. Hiding the boundary would manufacture confidence.
